%==============================================================================
% sections/02_literature_review.tex
%==============================================================================
\section{Literature Review and Positioning}
\label{sec:literature}

\subsection{Heavy-tailed sums and the single-big-jump principle}

The theoretical base is the large-deviation theory of subexponential and
regularly varying sums. In the finite-dimensional setting, the key heuristic is
that a high exceedance of a heavy-tailed sum is caused by one large term rather
than by a Gaussian accumulation of many small terms. This idea appears in the
standard treatments of regular variation and subexponentiality
\citep{BinghamGoldieTeugels1987,EmbrechtsKluppelbergMikosch1997,Resnick2007,FossKorshunovZachary2013}
and in random-walk large-deviation results such as
\citet{DenisovDiekerShneer2008}. The independent baseline of this paper is a
finite, non-identically distributed version of that principle, with the
additional goal of producing simulation estimators and financial risk forecasts.

\subsection{Rare-event simulation for heavy tails}

Conditional Monte Carlo for heavy-tailed sums was developed to avoid the poor
performance of crude Monte Carlo at very small probabilities.
\citet{AsmussenKroese2006} introduced an influential estimator based on
conditioning on all but a candidate large summand. Subsequent work studied
extensions, efficiency refinements, and related estimators, including
\citet{ChanKroese2011}, \citet{HartingerKortschak2009},
\citet{AsmussenKortschak2015}, and \citet{GhamamiRoss2012}. General simulation
criteria and rare-event efficiency concepts are reviewed in
\citet{AsmussenGlynn2007}, \citet{Glasserman2004}, and
\citet{JunejaShahabuddin2006}. The independent part of the present paper fits
this line but emphasizes non-identical tail constants, signed financial
contributions, exact active-set weights, finite-threshold concentration, and a
corrected finite-mean second-order expansion.

\subsection{Dependent heavy-tailed factor models}

The closest conceptual predecessor for dependent factor models is
\citet{DjehicheSvensson2007}, who study large deviations for sums generated by a
heavy-tailed factor structure and identify cases where factor and idiosyncratic
terms contribute to the tail probability. The earlier factor-model rare-event
work of \citet{PourbabaeeSolari2019} is the direct predecessor of the
independent baseline. The present draft bridges those two lines by making the
conditioning object explicit. Observed factor dependence can be handled by
moving from observed coordinates to latent shocks, independent blocks,
conditional copula kernels, or spectral directions.

\paragraph{Positioning relative to Pourbabaee--Solari (2019).}
\citet{PourbabaeeSolari2019} establish the first-order single-big-jump
equivalence and a CdMC estimator for an independent, non-identically distributed
finite-factor model. The present paper sharpens that result in four directions.
First, the active-set weights $c_i$ and the active-set constant $C$ enter the
limit theorem and the BRE bound symmetrically; the BRE bound $N^\alpha-1$ is
proved here in scale-invariant form (\cref{sec:independent-baseline}). Second,
the second-order expansion is corrected so that the local shift is supplied by
the leave-one-out mean $\mu_{-i}=\sum_{j\ne i}\E X_j$ rather than the full
portfolio mean; the resulting expansion passes an $N=1$ specialization check
that the prior expansion did not. Third, the dependent identity of
\cref{sec:dependent-cdmcs} extends the same partition argument to arbitrary
joint laws, recovering the independent case by setting the conditional kernel
equal to the marginal survival. Fourth, the multivariate-regular-variation and
hidden-regular-variation extensions of \cref{sec:mrv-spectral,sec:hidden-rv}
are new relative to the independent setting.

\subsection{Dependent rare-event simulation}

For dependent heavy-tailed random variables, the simulation literature includes
conditional Monte Carlo for log-elliptical structures and importance-sampling
methods under broader dependence. \citet{BlanchetRojasNandayapa2011} propose
asymptotically optimal algorithms for tail probabilities of sums of dependent
random variables, including a conditional Monte Carlo algorithm for elliptical
log variables and an importance sampler for general dependence under conditional
simulation assumptions. Recent copula-based work such as
\citet{ChengFuhPang2025} develops efficient importance sampling for copula
models. The present paper differs by keeping the factor-portfolio structure, the
regularly varying signed-tail constants, and the real-data VaR/ES workflow at
the center. Its dependent kernel estimator is exact whenever the conditional
survival is exact, while the shock, block, and spectral estimators choose more
stable conditioning objects when coordinate kernels are inefficient.

\subsection{Multivariate regular variation and hidden regular variation}

Multivariate regular variation is the natural language for dependent extremes.
\citet{HultLindskogMikoschSamorodnitsky2005} extend heavy-tailed large
deviations to multivariate regularly varying random walks and make precise the
idea that a single large vector step drives extremal behavior. The spectral
measure encodes whether extremes concentrate on coordinate axes, rays, faces, or
larger angular regions. \citet{SamorodnitskySun2016} develop multivariate
subexponential distributions and applications to ruin problems beyond the
simplest independent marginal setting.

Hidden regular variation is relevant when ordinary MRV is asymptotically
independent but joint-tail regions remain practically important.
\citet{MaulikResnick2004} characterize hidden regular variation and relate it to
residual tail dependence ideas. \citet{DasMitraResnick2013} emphasize that
classical MRV can understate probabilities of joint tail regions and that hidden
regular variation can improve estimates on such risk sets. The present paper
uses hidden regular variation as a second-order dependence correction layered on
top of the original marginal second-order expansion.

\subsection{Extreme-value dependence and financial tails}

Empirical financial returns are heavy-tailed, clustered, and dependent. The
stylized facts literature includes \citet{Cont2001} and the power-law review of
\citet{Gabaix2009}. Factor models and portfolio construction are grounded in
\citet{FamaFrench1993} and \citet{FamaFrench2015}. Quantitative risk management
and EVT-based VaR/ES methods are treated in \citet{McNeilFreyEmbrechts2015},
while simulation for heavy-tailed portfolio VaR is studied by
\citet{GlassermanHeidelbergerShahabuddin2002}. Backtesting methods include
\citet{Kupiec1995}, \citet{Christoffersen1998}, \citet{EngleManganelli2004},
and \citet{AcerbiSzekely2014}. Tail-dependence diagnostics build on
\citet{LedfordTawn1996} and \citet{ColesHeffernanTawn1999}.

\subsection{Recent developments and outstanding gaps}
\label{subsec:recent-developments}

Beyond the references above, the post-2020 literature on dependent
heavy-tailed simulation has continued to develop along several lines. Efficient
importance-sampling schemes for copula and vine-copula models
\citep{ChengFuhPang2025} broaden the class of dependence structures for which
provably efficient estimators exist. Work on multivariate regular variation
under partial mass on cones, on hidden regular variation with multiple
asymptotic-independence scales, and on tail-process methods for serially
dependent extremes (in the line of \citet{BasrakSegers2009}) refines the
extremal-dependence apparatus. Empirical work on tail-index stability over
crisis and non-crisis subperiods continues to motivate rolling-window
diagnostics. The present paper integrates these strands at the level of the
factor-portfolio model: it states the dependent CdMC identity in a form that
specializes to each of latent-shock, block, copula, MRV, and HRV structures,
and supplies a real-data protocol that reports the diagnostics needed to
choose among them.

\paragraph{Outstanding gaps addressed here.}
Three gaps motivated the present paper. First, the existing CdMC literature
gives BRE results for restricted classes of dependent sums; a unified
statement that recovers the independent BRE bound $N^\alpha-1$ as a special
case has been absent. Second, the second-order expansion for finite-portfolio
heavy-tailed sums was, prior to the present paper, vulnerable to a
specialization failure at $N=1$ (the local-shift correction term should
vanish when there is only one summand). Third, the empirical literature on
factor-portfolio tail risk has historically relied on either independent
parametric models or fully nonparametric extreme-value methods; an
intermediate class that exploits factor structure while admitting principled
dependence diagnostics has been missing.
\Cref{sec:independent-baseline,sec:dependent-cdmcs,sec:mrv-spectral,sec:hidden-rv,sec:estimators-efficiency,sec:real-data}
address these gaps in order.

\begin{table}[p]
\centering
\begin{tabularx}{\linewidth}{p{0.19\linewidth}p{0.24\linewidth}p{0.22\linewidth}X}
\toprule
Literature & Representative references & What is established & Role of this paper \\
\midrule
Independent heavy-tailed CdMC & \citet{AsmussenKroese2006}, \citet{ChanKroese2011}, \citet{HartingerKortschak2009}, \citet{AsmussenKortschak2015} & Conditional Monte Carlo and efficiency for heavy-tailed sums & Extends i.n.i.d. finite-factor setup, corrected second order, financial pipeline \\
Dependent heavy-tailed factor models & \citet{DjehicheSvensson2007}, \citet{PourbabaeeSolari2019} & Factor/idiosyncratic large deviations and simulation ideas & Adds shock-basis constants, estimator hierarchy, real-data diagnostics \\
Dependent rare-event simulation & \citet{BlanchetRojasNandayapa2011}, \citet{ChengFuhPang2025} & Efficient simulation for log-elliptical, copula, and broad dependent structures & Specializes to regularly varying factor portfolios with active sets and spectral controls \\
MRV and hidden RV & \citet{HultLindskogMikoschSamorodnitsky2005}, \citet{MaulikResnick2004}, \citet{DasMitraResnick2013}, \citet{SamorodnitskySun2016} & Vector-tail measures, one large vector step, hidden joint-tail regions & Converts theory into CdMC estimators and financial diagnostics \\
Financial tail risk & \citet{GlassermanHeidelbergerShahabuddin2002}, \citet{McNeilFreyEmbrechts2015}, \citet{FamaFrench2015} & Heavy-tailed VaR/ES modeling and asset-pricing factors & Provides dependent-factor VaR/ES protocol with backtests and attribution \\
\bottomrule
\end{tabularx}

\caption[Literature map]{Literature map and positioning of the expanded paper.
The contribution column states the additional role of the present draft beyond
the cited line of work.}
\label{tab:literature-map}
\end{table}

\begin{table}[p]
\centering
\begin{tabularx}{\linewidth}{p{0.04\linewidth}p{0.27\linewidth}p{0.13\linewidth}p{0.14\linewidth}p{0.13\linewidth}X}
\toprule
Rank & Module & Innovation & Contribution & Feasibility & Recommended status \\
\midrule
1 & MRV spectral/radial CdMC & 5/5 & 5/5 & 3/5 & Theoretical centerpiece; best generalization of one big jump \\
2 & Hidden-cone second-order correction & 5/5 & 4.5/5 & 2/5 & High-risk/high-reward; include as template and simulation target \\
3 & Latent/common-shock factors & 3.5/5 & 4.5/5 & 4.5/5 & Best near-term extension and easiest financial interpretation \\
4 & Copula/vine conditional CdMC & 4/5 & 4/5 & 3.5/5 & Strong applied estimator when conditional kernels are reliable \\
5 & Block-dependent CdMC & 3/5 & 4/5 & 4/5 & Useful empirical bridge and robust implementation option \\
6 & Serial tail-process extension & 4/5 & 3.5/5 & 2.5/5 & Later multi-period extension; not central to this draft \\
\bottomrule
\end{tabularx}

\caption[Extension ranking]{Ranking of extension modules by innovation,
expected contribution, feasibility, and recommended development status. Scores
are qualitative design scores on a five-point scale, not empirical results.}
\label{tab:extension-ranking}
\end{table}
